\chapter{Biological and Epigenetic Background}
\label{chap:biological_background}

% -------------------------------------------------------
\section{Breast Cancer: Epidemiology and Molecular Subtypes}
% -------------------------------------------------------

Breast cancer is the most frequently diagnosed malignancy in women worldwide and
represents a major public health burden across both high- and low-income countries
\cite{Hanahan2011}. Its clinical and biological heterogeneity is one of the most
distinctive features of the disease: not all breast cancers behave alike, progress
at the same rate, or respond to the same treatments.
At the histological level, the vast majority of breast malignancies originate from
the epithelial compartment of the mammary gland, most commonly from the luminal
cells lining the ductal and lobular structures. The classical progression model
describes a continuum from normal epithelium through non-atypical hyperplasia,
atypical ductal hyperplasia (ADH), ductal carcinoma in situ (DCIS), and ultimately
invasive carcinoma \cite{Hanahan2000}:
\[
\text{Normal epithelium} \rightarrow
\text{Hyperplasia} \rightarrow
\text{DCIS} \rightarrow
\text{Invasive carcinoma}.
\]
At the molecular level, tumors are stratified according to the expression status of
estrogen receptor (ER), progesterone receptor (PR), and HER2, defining biologically
distinct subtypes — luminal A, luminal B, HER2-enriched, and triple-negative — with
different prognostic and therapeutic implications. This molecular classification
reflects underlying differences in the cell of origin, epigenetic landscape, and
selective pressures operative during tumor evolution.
A subset of breast cancers arises on an inherited background of germline susceptibility.
Pathogenic variants in \textit{BRCA1} and \textit{BRCA2}, genes encoding key
components of the homologous recombination DNA repair machinery, substantially
increase lifetime risk \cite{NCI_BRCA}. BRCA-associated tumors tend to be
triple-negative and exhibit distinct baseline epigenetic profiles that differ
fundamentally from sporadic disease \cite{Suijkerbuijk2008,Bartlett2022}.
For this reason, BRCA mutation carriers require separate analytical treatment in
any study of epigenetic field effects, as discussed in Chapter~\ref{sec:data_preparation_exploration}.

% -------------------------------------------------------
\section{The Hallmarks Framework and Epigenetic Reprogramming}
% -------------------------------------------------------

The conceptual framework of cancer biology has been progressively organized
around the notion of ``hallmarks'': a set of functional capabilities that normal
cells must acquire to become malignant \cite{Hanahan2011, Hanahan2000}. These
include sustained proliferative signaling, resistance to growth suppression, evasion
of apoptosis, unlimited replicative potential, induction of angiogenesis, and the
activation of invasion and metastasis.
Importantly, the updated hallmarks framework explicitly recognizes epigenetic
reprogramming as a core enabling mechanism \cite{Hanahan2011}. Tumor cells do not
merely accumulate genetic mutations; they undergo widespread reorganization of
their chromatin architecture and DNA methylation landscape, which silences
tumor-suppressive programs and derepresses oncogenic pathways. This repositions
epigenetic alterations not as secondary consequences of malignant transformation
but as primary and early drivers of the tumorigenic process.

% -------------------------------------------------------
\section{DNA Methylation: Mechanism and Genomic Architecture}
\label{sec:methylation_mechanism}
% -------------------------------------------------------

DNA methylation consists of the covalent addition of a methyl group to the
5-carbon position of cytosine residues within CpG dinucleotides, a reaction
catalyzed by a family of DNA methyltransferases (DNMT3A, DNMT3B for \textit{de novo}
methylation; DNMT1 for maintenance after replication) \cite{bird2002}.
CpG dinucleotides are non-randomly distributed in the genome: statistically
expected at roughly one occurrence per 16 bp, they are in practice four- to
fivefold depleted genome-wide — a depletion attributed to the historical
mutability of methylated cytosines, which deaminate spontaneously to thymine.
The exception to this depletion is found at \textit{CpG islands}: dense,
GC-rich clusters of CpG dinucleotides, typically 500–2000 bp in length,
located at the promoters of approximately 60\% of human protein-coding genes
\cite{bird2002}. In normal somatic cells, CpG islands are kept constitutively
unmethylated by active mechanisms, preserving transcriptional access.

\subsection{CpG Islands, Shores, Shelves, and Open Sea}

The regulatory significance of DNA methylation varies substantially depending
on where in the genome a CpG locus resides. A widely adopted classification,
used in Illumina array annotation, distinguishes four positional contexts
relative to CpG islands \cite{Irizarry2009}.
\begin{itemize}[label=--]
    \item \textbf{CpG Islands}: regions with high CpG density and high
    GC content, operationally defined as stretches of $\geq 200$ bp with
    CpG observed/expected ratio $> 0.6$ and GC content $> 50\%$.
    Methylation at island-overlapping promoters is the canonical mechanism
    of transcriptional silencing in both development and cancer.

    \item \textbf{Shores}: regions within 2 kb flanking a CpG island on
    either side. Although less CpG-dense than islands, shores are the sites
    of the most dynamic and biologically informative methylation variation.
    Landmark work by Irizarry et al.\ demonstrated that tissue-specific
    methylation differences between normal cell types are concentrated at
    shores rather than at islands, and that cancer-associated hypermethylation
    also preferentially targets shores \cite{Irizarry2009}. This makes shores
    the primary locus of epigenetic switching between cell states.

    \item \textbf{Shelves}: regions 2–4 kb from the island boundary. Their
    methylation patterns are less variable than shores and less directly tied
    to promoter regulation, though they retain some contextual regulatory
    relevance.

    \item \textbf{Open Sea}: CpG loci not associated with any island, shore,
    or shelf. These constitute the majority of CpGs in the genome ($\sim$70\%)
    and tend to be constitutively methylated in somatic tissues. Their
    methylation changes are more often associated with global hypomethylation
    and genomic instability than with specific gene regulation.
\end{itemize}

This spatial hierarchy has direct consequences for how methylation data
should be interpreted and weighted in analytical frameworks. Loci overlapping
islands and shores are enriched for functional regulatory activity — they are
more likely to be associated with gene expression changes when their methylation
state shifts. Shelf and open-sea loci, by contrast, tend to reflect structural
genomic states rather than specific transcriptional events. When constructing
feature selection or scoring systems based on biological relevance, this gradient
from island to open sea provides a principled basis for differential weighting.

\subsection{Genomic Position Determines the Functional Consequence of Methylation}

A further distinction of biological relevance concerns the relationship between
methylation and transcription depending on the position of the CpG locus relative
to the gene structure. The classical rule — methylation silences expression — applies
specifically to \textit{promoter-associated} CpG islands. Within the gene body
(introns and exons), the relationship is inverted: gene body methylation correlates
\textit{positively} with transcriptional activity \cite{Jones2012}. This is thought
to reflect the need to suppress spurious transcription initiation within highly
transcribed genes. The functional interpretation of a methylation change therefore
depends critically on whether the affected locus falls in a promoter or an
intragenic context.
In genome-wide methylation profiling studies, the Illumina array annotation
provides positional information (TSS200, TSS1500, 5'UTR, 1stExon, Body, 3'UTR)
that, combined with the island/shore/shelf classification, allows a nuanced
characterization of each locus. CpGs annotated to TSS regions within islands or
shores carry the highest prior probability of regulatory relevance and are therefore
the most informative for studies of transcriptional silencing in cancer.

% -------------------------------------------------------
\section{Epigenetic Alterations in Cancer: Too Much and Too Little}
% -------------------------------------------------------

Cancer genomes are characterized by a paradoxical co-occurrence of global
hypomethylation and focal hypermethylation \cite{ehrlich2002}, a duality
confirmed across virtually all tumor types, including breast cancer.
\emph{Global hypomethylation} affects repetitive sequences (satellite DNA,
LINE-1 elements) and promotes genomic instability through aberrant activation
of transposable elements and loss of imprinting, resulting in transcriptional
derepression at normally silenced genes \cite{ehrlich2002}.
\emph{Focal hypermethylation}, by contrast, targets CpG islands at the promoters
of specific genes, leading to heritable transcriptional silencing. Affected genes
frequently include classical tumor suppressors (e.g., \textit{CDKN2A},
\textit{BRCA1}, \textit{MLH1}), as well as genes involved in apoptosis, cell
adhesion, and developmental signaling. This epigenetic silencing is functionally
equivalent to mutation but potentially reversible, making it an attractive
therapeutic target.
Critically, neither phenomenon is random. Genes susceptible to hypermethylation
in cancer are systematically enriched for targets of the Polycomb Repressive
Complex 2 (PRC2) and display bivalent chromatin states in embryonic stem cells ---
a configuration marking developmental genes as both transcriptionally poised and
repressed \cite{Teschendorff2010_stemcell}, suggesting that epigenetic alterations
in malignancy partially recapitulate stem cell programs.
% -------------------------------------------------------
\section{Epigenetic Drift, Aging, and Cancer Risk}
% -------------------------------------------------------

Aging is the single strongest risk factor for breast cancer. The incidence of most
breast malignancies rises sharply with age, and this epidemiological observation
has a direct molecular correlate in the epigenome.
Over the course of a lifetime, DNA methylation patterns undergo progressive and
partially stochastic changes collectively referred to as \textit{epigenetic drift}
\cite{Horvath2013}. This phenomenon is characterized by gradual increases in
methylation variability across both CpG loci and individuals. The drift is not
uniform: it preferentially affects genes associated with Polycomb repression and
bivalent chromatin — precisely the same loci that tend to become hypermethylated
in cancer \cite{Teschendorff2010,Teschendorff2010_stemcell}.
Horvath's landmark study demonstrated that DNA methylation patterns can be
exploited to construct a highly accurate molecular clock: the ``epigenetic age''
of a tissue, estimated from a weighted combination of CpG methylation values,
correlates strongly with chronological age across more than 50 tissue types
\cite{Horvath2013}. Crucially, epigenetic age acceleration — a discordance between
molecular and chronological age — has been associated with cancer risk, suggesting
that the pace of epigenetic drift itself may be a biological marker of malignant
susceptibility.
From a mathematical perspective, aging-associated drift can be modeled as a
diffusion-like process in epigenetic state space: methylation levels at susceptible
loci undergo a gradual increase in inter-individual variance over time. This growing
dispersion may create a permissive landscape for malignant transformation by
producing a fraction of cells with aberrant epigenetic configurations from which
selection can act.
An important study by Teschendorff et al.\ in breast tissue documented precisely
this mechanism: aging-related DNA methylation changes in normal breast tissue
were found to target the same genomic regions altered in cancer, establishing a
quantitative link between physiological aging and pre-malignant epigenetic
reprogramming \cite{Teschendorff2019}.


% -------------------------------------------------------
\section{Field Cancerization: The Concept and Its Epigenetic Basis}
% -------------------------------------------------------

A central concept motivating the present thesis is that of \textit{field
cancerization}. Originally proposed by Slaughter in 1953 from histopathological
observations in oral cancer, the field effect hypothesis posits that
carcinogenesis does not occur in isolation but within a ``conditioned'' tissue
field broadly reprogrammed by molecular alterations preceding frank malignancy.
In breast cancer, epigenetic field defects have been systematically characterized
using genome-wide methylation profiling. Teschendorff et al.\ demonstrated that
histologically normal breast tissue adjacent to tumors harbors a specific pattern
of methylation alterations distinct from both distant normal tissue and tumor
tissue \cite{Teschendorff2016}, enriched at PRC2 target genes, CTCF-binding
regions, and loci involved in WNT and FGF developmental pathways — precisely
the categories implicated in the aging-cancer methylation axis described above.
Crucially, this field defect signature was detectable not through classical
differential mean analysis but through \textit{differential variability}, a
framework quantifying the increase in methylation dispersion across samples
rather than directional shifts in average levels \cite{Teschendorff2016,Teschendorff2014}.
Under a classical framework, loci with increased variance but small mean
differences are invisible to standard hypothesis tests and require
variance-sensitive statistics such as the Levene test or its robust variants.
The biological interpretation is that field defects reflect a partial and
heterogeneous epigenetic reprogramming, where the stochastic nature of the
alteration manifests as increased dispersion across individuals rather than
a uniform shift.
The progressive nature of this process is further supported by the observation
that hypervariable loci in adjacent tissue tend to become hypermethylated in
the matched invasive tumor, suggesting a deterministic convergence from a
heterogeneous pre-malignant landscape to a consolidated malignant state
\cite{Teschendorff2016}:
\[
\text{Normal} \rightarrow
\text{Epigenetically reprogrammed adjacent tissue} \rightarrow
\text{Invasive carcinoma}.
\]
This model implies that adjacent tissue represents an earlier stage on the
same carcinogenic trajectory rather than mere contamination from the nearby
tumor, making the detection and characterization of its methylation signature
both biologically and clinically significant.

% -------------------------------------------------------
\section{Differential Variability in Cancer}
% -------------------------------------------------------

The discovery that field defects are better captured by differential variability
than by differential means represents a genuine shift in the analytical paradigm
for epigenetic cancer studies \cite{Teschendorff2014}.
Formally, let $X_j$ denote the methylation level at locus $j$ across a cohort
of samples. Classical differential methylation analysis focuses on the comparison
of group means:
\begin{equation}
\Delta_j = \mathbb{E}[X_j \mid A] - \mathbb{E}[X_j \mid B].
\end{equation}
Differential variability analysis instead targets the comparison of group
variances:
\begin{equation}
\delta\sigma^2_j = \text{Var}(X_j \mid A) - \text{Var}(X_j \mid B).
\end{equation}
A locus is said to be \textit{differentially variable} (DV) if $\delta\sigma^2_j$
is significantly non-zero. Under a pre-malignant model, one expects $\delta\sigma^2_j > 0$
in normal-adjacent tissue relative to distant normal tissue, reflecting increased
epigenetic disorder in the field-affected region.
Teschendorff et al.\ formalized this framework and developed statistical tools
for its application to genome-wide methylation data \cite{Teschendorff2014}.
Their analysis demonstrated that DV CpGs in normal-adjacent breast tissue are
substantially enriched for biologically meaningful genomic features (PRC2 targets,
bivalent domains) and display stronger associations with cancer risk markers —
including tumor proliferation as measured by KI67 — than differentially methylated
CpGs identified by mean-based tests alone \cite{Teschendorff2016}.
The implication for any analytical framework aimed at characterizing early
epigenetic alterations is direct: a reliance on mean-based feature selection alone
will systematically miss a biologically relevant component of the pre-malignant
signal. This motivates the integration of variability-aware preprocessing and
feature selection strategies, which are developed in the methodological chapters
of this thesis.

\section{The Normal--Adjacent--Tumor Axis as an Analytical Framework}
\label{sec:bio_synthesis}

Breast carcinogenesis is not a sudden event but a progressive deviation from
epigenetic homeostasis, driven by the interplay of aging, stochastic drift, and
selective pressures across a population of cells. The biological evidence reviewed
in this chapter converges on a model in which this progression is both measurable
and partially predictable from methylation data alone.
Normal breast tissue accumulates methylation variability with age at loci
associated with Polycomb regulation and bivalent chromatin, creating a
heterogeneous epigenetic landscape from which pre-malignant field defects emerge.
Adjacent tissue, histologically indistinguishable from normal, harbors a
measurable epigenetic signature of this reprogramming — detectable predominantly
as increased variability rather than a directional mean shift. As carcinogenesis
proceeds, this heterogeneous pre-malignant state converges toward a consolidated
malignant configuration characterized by focal hypermethylation at
tumor-suppressive loci and global hypomethylation elsewhere.
This progressive model has three direct consequences for the quantitative analyses
developed in this thesis. First, the Normal--Adjacent--Tumor axis should be treated
as a continuous biological spectrum: adjacent tissue occupies a genuinely
intermediate epigenetic state, distinct from both normal and tumor, and warrants
explicit characterization rather than collapsing into either extreme. Second,
feature selection must be sensitive to variance differences and not only to mean
shifts, as field defects are heterogeneous signals that may not survive standard
mean-based filtering. Third, cross-cohort reproducibility is a necessary
validation criterion: signatures replicating across independent datasets are more
likely to reflect genuine biological structure than cohort-specific noise.
These three principles — spectrum-aware design, variance sensitivity, and
cross-cohort validation — directly motivate the construction, preprocessing, and
comparative characterization of the three cohorts introduced in
Chapter~\ref{sec:data_preparation_exploration}.